\begin{frame}[t]{跨模态 CPT：蛋白组语料与 RNA 对照}
\fontsize{8.8}{12}\selectfont\renewcommand{\arraystretch}{1.14}
\begin{tabularx}{\textwidth}{@{}P{2.5cm}Y Y@{}}
\toprule\textbf{来源}&\textbf{数据组成与大小}&\textbf{处理和用途}\\\midrule
公开蛋白组语料 & 440 来源 / 48,843 样本；细胞系 21,910，组织 15,023，体液 6,578，其他 5,332 & 逐来源解析、AI agent 辅助 + 专家复核；部分原始谱图重搜；不足 10 样本的来源被排除\\
定量策略 & TMT/iTRAQ 15,526；DIA 14,170；LFQ 12,541；其他 6,606 & 平均深度 4,223 蛋白/样本；不同技术与样本类型分布不均\\
ARCHS4 human v2.5 & 888,821 原始 bulk RNA 样本中，877,715 成功编码 & 额外 RNA-only 对照，以及 RNA + 蛋白混合训练\\
\bottomrule
\end{tabularx}
\vspace{.22cm}
\subhead{三组对照：仅蛋白组 / 仅 bulk RNA / 蛋白组 + bulk RNA}
\remark{均报告 1 epoch；相同 epoch 不代表相同总训练步数。正文另写 RNA 训练有 1,003,000 行，和 877,715 唯一样本如何对应未解释，复现时需核实。}
\source{[25] §2.1--2.2、Fig. 1、SI Table S1；这里的“其他样本类型”包括原代、干细胞/类器官与标签不明来源，不等于“其他定量策略”。}
\end{frame}
